%!TeX root=MemoriaTFG.tex

\chapter{Materiales y métodos}
\label{cap:metodos}

\section{Datasets}
\label{sec:datos}

La evaluación se realiza sobre doce datasets externos públicos
agrupados por modalidad y finalidad. La
tabla~\ref{tab:datasets-resumen} sintetiza sus características.

\begin{table}[H]
\centering
\footnotesize
\setlength{\tabcolsep}{4pt}
\renewcommand{\arraystretch}{0.95}
\caption{Datasets externos empleados en la evaluación.
$N_{\mathrm{test}}$: tamaño nominal del \emph{split} de evaluación
canónico; ``Clases'': número de categorías diagnósticas o conceptos
en la convención propia del dataset. El asterisco
(\textsuperscript{$\ast$}) marca los datasets con solapamiento
documentado frente al dataset de preentrenamiento Derm1M
(auditoría en secci\'on~\ref{sec:auditoria}). Las dos variantes de Derm7pt
(clínica y dermatoscópica) comparten origen
(Seven-Point Checklist) y se reportan como un único dataset en el
conteo agregado de doce datasets externos.}
\label{tab:datasets-resumen}
\begin{tabular}{lrcll}
\hline
Dataset & $N_{\mathrm{test}}$ & Clases & Modalidad & Uso \\
\hline
HAM10000~\cite{ham10000}             & 1\,232  & 7        & Dermoscopia      & LP, FT, ZS, LLMs \\
BCN20000~\cite{bcn20000}             & 1\,242  & 9        & Dermoscopia      & LP \\
PAD-UFES-20~\cite{padufes}           & 461     & 6        & Clínica móvil    & LP, FT \\
DDI~\cite{ddi}                       & 137     & 2        & Clínica          & LP, FT, equidad \\
Dermnet~\cite{dermnet}               & 4\,002  & 23       & Clínica          & LP\textsuperscript{$\ast$} \\
Derm7pt clínico~\cite{kawahara2019}  & 168     & 2        & Clínica          & LP \\
Derm7pt dermo~\cite{kawahara2019}    & 225     & 2        & Dermoscopia      & LP \\
HIBA                                 & 334     & 2        & Dermoscopia      & LP\textsuperscript{$\ast$} \\
MSKCC                                & 1\,664  & 2        & Dermoscopia      & LP, FT\textsuperscript{$\ast$} \\
WSI patches                          & 12\,354 & 16       & Histopatología   & LP \\
ISIC2018~\cite{isic2018}             & 2\,594  & binaria  & Dermoscopia      & Segmentación \\
Fitzpatrick17k~\cite{fitzpatrick17k} & 16\,577 & 114 + FP & Clínica          & Equidad \\
SkinCon~\cite{skincon}               & 3\,230  & 48 conc. & Clínica          & CBM \\
\hline
\end{tabular}
\end{table}

Los conjuntos dermatoscópicos (HAM10000, BCN20000, Derm7pt dermo,
HIBA, MSKCC, ISIC2018) cubren la modalidad de imagen polarizada
adquirida con dermatoscopio. Los conjuntos de fotografía clínica
(PAD-UFES-20, DDI, Dermnet, Derm7pt clínico) reproducen las
condiciones de captura en consulta y atención primaria. WSI patches
recoge $12\,354$ \emph{parches} histopatológicos en $16$
categorías. Fitzpatrick17k introduce el eje de equidad mediante
anotación conjunta de patología ($114$ categorías) y fototipo
cutáneo (escala Fitzpatrick I-VI). SkinCon proporciona anotación
densa basada en $48$ conceptos clínicos visuales para la
evaluación de Concept Bottleneck Models. Las particiones
\emph{train/val/test} utilizadas son las canónicas declaradas por
los responsables de cada dataset, sin reasignación.

\section{Ontología jerárquica L1/L2/L3}
\label{sec:ontologia}

La heterogeneidad de las taxonomías nativas impide la comparación
directa entre datasets. Para resolverla hemos diseñado, en
colaboración con la Dra.~R.~Taberner, una ontología jerárquica de
tres niveles que sirve de vocabulario común para todos los protocolos
del trabajo. El nivel L1 (etiológico) agrupa los diagnósticos
en cuatro familias: patología inflamatoria, tumoral, infecciosa y
genodermatosis. El nivel L2 (subcategoría clínica) contiene
$43$ entradas; cinco de ellas no se desagregan al nivel L3 al
considerarse suficientemente específicas (\emph{Esclerosis
tuberosa}, \emph{Neurofibromatosis tipo~1}, \emph{Incontinentia
pigmenti}, \emph{Queratodermia palmoplantar hereditaria} y
\emph{Pseudoxantoma elástico}). El nivel L3 (diagnóstico) contiene
$367$ entradas individuales. Cada entrada L3 lleva asociados
códigos SNOMED~CT y CIE-10 para garantizar interoperabilidad con
sistemas clínicos. Once de los doce datasets se mapean a la
ontología; SkinCon queda fuera al estar anotado con conceptos en
lugar de diagnósticos. El dataset armonizado agregado contiene
$72\,654$ imágenes etiquetadas sobre vocabulario común.

\section{Dataset DermapixelAI 1.0}
\label{sec:dermapixel}

DermapixelAI~1.0 es un dataset dermatológico en español, derivado
del archivo del blog Dermapixel mantenido por la Dra.~R.~Taberner.
La versión utilizada contiene $1\,089$ imágenes vinculadas a
$669$ casos clínicos con texto narrativo asociado (mediana de
$223$ palabras por caso), mayoritariamente fotografía clínica
($1\,036$) con un subconjunto dermatoscópico reducido
($49$). Las imágenes se etiquetan conforme a la ontología
L1/L2/L3 y se particionan en \emph{train/val/test} bajo regla
\emph{case-aware} ($891 / 157 / 41$ imágenes; $538 / 100 / 31$
casos), garantizando que las imágenes de un mismo caso clínico
no se distribuyen entre splits.
El $97{,}89\%$ del dataset está mapeado a la ontología L1/L2/L3
con revisión experta del vocabulario
(\texttt{label\_source = ontology}); el $84{,}39\%$ ha pasado
además revisión visual caso-a-caso por la Dra.~R.~Taberner
(\texttt{rosa\_verified}). Las dos cifras no son intercambiables y
se documentan en detalle, junto a un tercer campo de validación
textual, en el Anexo~\ref{cap:anexoc}. El dataset se distribuye
bajo licencia CC-BY-NC-SA~4.0. El pipeline detallado de
construcción, caracterización completa y datasheet figura en el
Anexo~\ref{cap:anexoc}; el documento formal de entrega, en el
Anexo~\ref{cap:anexoe}. La explotación experimental sistemática
sobre DermapixelAI~1.0 se inicia hacia el cierre del periodo
cubierto por este documento; las cifras del
capítulo~\ref{cap:resultados} no incluyen evaluación sobre este
dataset.

\section{Modelos evaluados}
\label{sec:modelos}

La evaluación abarca catorce modelos agrupados en cuatro familias.
La tabla~\ref{tab:modelos} sintetiza sus características.

\begin{table}[H]
\centering
\footnotesize
\setlength{\tabcolsep}{4pt}
\renewcommand{\arraystretch}{0.95}
\caption{Modelos evaluados.}
\label{tab:modelos}
\begin{tabular}{llrl}
\hline
Modelo & Familia & Parámetros & Referencia \\
\hline
PanDerm Base                & Encoder visual derm.       & $\sim 86$\,M & \cite{panderm2025} \\
PanDerm Large               & Encoder visual derm.       & $\sim 307$\,M & \cite{panderm2025} \\
DINOv2 ViT-L/14             & Encoder visual general     & $\sim 300$\,M & \cite{dinov2} \\
ConvNeXt-Large              & ConvNet supervisada        & $\sim 198$\,M & --- \\
EfficientNetV2-Large        & ConvNet supervisada        & $\sim 118$\,M & --- \\
DermLIP v1 (PubMedBERT)     & Vision-lenguaje derm.      & --- & \cite{derm1m} \\
DermLIP v2 (PubMedBERT)     & Vision-lenguaje derm.      & --- & \cite{derm1m} \\
DermLIP original (GPT-2)    & Vision-lenguaje derm.      & --- & \cite{derm1m} \\
BiomedCLIP                  & Vision-lenguaje biomédico  & --- & \cite{biomedclip} \\
SigLIP-Large SO400M         & Vision-lenguaje general    & $\sim 878$\,M & \cite{siglip2023} \\
SAM2.1-Large + LoRA         & Segmentación               & $\sim 227$\,M & \cite{sam2} \\
GPT-4o                      & LLM multimodal             & no público & --- \\
Gemini 2.5 Pro              & LLM multimodal             & no público & --- \\
MedGemma 4B Vision          & LLM multimodal biomédico   & $\sim 4$\,B & --- \\
MedGemma 27B Text+LoRA      & LLM multimodal biomédico   & $\sim 27$\,B & --- \\
BLIP-2 (Q-Former)           & Vision-lenguaje general    & --- & \cite{blip2_2023} \\
\hline
\end{tabular}
\end{table}

PanDerm Base produce \emph{embeddings} de $768$ dimensiones;
PanDerm Large y DINOv2 ViT-L/14, $1\,024$ dimensiones; SigLIP-Large
SO400M, $1\,152$ dimensiones. La cifra de $\sim 307$\,M parámetros
indicada para PanDerm Large corresponde al valor nominal del
codificador ViT-L/16 publicado por los autores; al cargarlo como
extractor de características y sustituir la cabeza original por la
identidad se obtienen $303{,}3$\,M parámetros efectivos, diferencia
sin impacto sobre el rendimiento al no participar la cabeza
descartada en la inferencia. La familia DermLIP emplea ViT-B/16
como encoder visual y dos encoders textuales alternativos:
PubMedBERT extendido a $256$ tokens (variantes v1 y v2) y GPT-2
con tokenizador de CLIP a $77$ tokens (variante original).
DermFM-Zero~\cite{dermfmzero2026} no se evalúa empíricamente al no
estar disponibles sus pesos en la fecha de cierre.

\section{Infraestructura experimental}
\label{sec:infraestructura}

Los experimentos se ejecutan sobre una NVIDIA DGX Spark equipada
con chip Grace Hopper GB10, $128$\,GB de memoria unificada
compartida entre CPU y GPU, sistema Linux ARM aarch64 y CUDA~13.0.
La precisión nativa es BF16 (\emph{brain floating-point} de 16
bits) en inferencia y entrenamiento, con FP32 en optimizador y
pérdida en los protocolos de ajuste fino (precisión mixta
canónica). La memoria unificada permite mantener cargados modelos
del orden de $55$\,GB en BF16 sin \emph{offloading} ni
cuantización; en particular, MedGemma~27B opera localmente sin
compresión. El entorno de software está fijado en Python~3.12.3,
PyTorch~2.11.0, \texttt{timm}~0.9.16, \texttt{scikit-learn},
\texttt{transformers}, \texttt{open\_clip}, \texttt{sam2},
\texttt{peft} y FAISS~\cite{faiss2019}. Los tiempos de cálculo
característicos sobre este equipo son los siguientes: sondeo lineal
completo de un encoder sobre un dataset estándar, $1$--$3$
minutos; ajuste fino de PanDerm Large sobre HAM10000 a $50$
épocas con TTA, $\sim 15$ horas; ajuste fino de SAM2.1-Large
sobre ISIC2018 con LoRA a $50$ épocas, $\sim 24$ horas;
entrenamiento de un Sparse Autoencoder Large ($16\,384$
características) sobre los \emph{embeddings} de PanDerm,
$\sim 4$ horas.

\section{Protocolos experimentales}
\label{sec:protocolos}

\subsection{Sondeo lineal (LP)}
\label{sec:lp}

El protocolo congela el encoder visual y entrena un clasificador
lineal sobre los \emph{embeddings} extraídos. La extracción emplea
la normalización canónica de cada modelo: ImageNet
($\mu=[0{,}485, 0{,}456, 0{,}406]$,
$\sigma=[0{,}229, 0{,}224, 0{,}225]$) para PanDerm y DINOv2;
parámetros OpenAI CLIP para la familia DermLIP, BiomedCLIP y
SigLIP. Los vectores resultantes se almacenan en disco para evitar
recomputaciones entre experimentos. El clasificador es una
regresión logística multinomial con optimizador L-BFGS, penalización
L2 con coeficiente $C = 1{,}0$, \texttt{max\_iter}~$=5\,000$ y
\texttt{random\_state}~$=42$. Los \emph{splits} son las particiones
canónicas declaradas en la tabla~\ref{tab:datasets-resumen}. El
protocolo se aplica sobre PanDerm Base y Large, DINOv2,
ConvNeXt-Large, EfficientNetV2-Large y la rama visual de DermLIP v2,
BiomedCLIP y SigLIP-Large.

\subsection{Ajuste fino supervisado (FT)}
\label{sec:ft}

El protocolo ajusta conjuntamente el encoder visual y una cabeza de
clasificación específica. La configuración reproduce la del paper
original de PanDerm~\cite{panderm2025}: optimizador AdamW con
\emph{weight decay} $0{,}05$; tasa de aprendizaje inicial
$\eta = 5\times 10^{-4}$; \emph{cosine annealing} con
\emph{warmup} lineal de $10$ épocas; \emph{layer decay}
$0{,}65$; \emph{drop-path} $0{,}2$;
\emph{label smoothing} $\epsilon = 0{,}1$;
Mixup $\alpha = 0{,}8$; CutMix $\alpha = 1{,}0$;
\emph{weighted random sampler} con pesos inversamente proporcionales
a la frecuencia de cada clase; $50$ épocas; precisión mixta
canónica (BF16 en el cuerpo del modelo, FP32 en optimizador y
pérdida); semillas de \texttt{torch}, \texttt{numpy} y
\texttt{random} fijadas a $0$. En la evaluación final sobre
HAM10000 se aplica \emph{test-time augmentation} (TTA) promediando
los logits sobre $5$ augmentaciones deterministas (volteos
horizontal y vertical, rotación, \emph{jitter} de color) por
imagen. El protocolo se aplica sobre PanDerm Base y Large en
HAM10000, PAD-UFES-20, DDI y MSKCC.

\subsection{Eficiencia en etiquetas}
\label{sec:label-eff}

El protocolo mide la dependencia del rendimiento respecto al
volumen de datos etiquetados disponibles, manteniendo congelado el
encoder. Sobre HAM10000 y PAD-UFES-20 se entrenan clasificadores
lineales (con la misma configuración del sondeo lineal,
secci\'on~\ref{sec:lp}) con submuestreos estratificados del conjunto de
entrenamiento al $1\%$, $5\%$, $10\%$, $20\%$,
$50\%$ y $100\%$. La estratificación preserva la
proporción original de clases en cada subconjunto. Los submuestreos
se generan con \texttt{random\_state}~$=42$ para garantizar
reproducibilidad. PanDerm Large es el encoder principal; PanDerm
Base se utiliza como referencia adicional sobre HAM10000.

\subsection{Segmentación binaria de lesión}
\label{sec:segmentacion}

SAM2.1-Large se adapta al dominio dermatológico mediante
\emph{fine-tuning} con LoRA~\cite{lora}. Las componentes ajustadas
son el decodificador de máscaras (\texttt{mask\_decoder}) y el
\emph{promptable encoder}; el \texttt{image\_encoder} (Hiera-Large
preentrenado) permanece congelado, lo que mantiene fijos la
mayoría de los parámetros del modelo. La configuración LoRA emplea
\emph{rank} $r = 8$ con la implementación estándar de la biblioteca
\texttt{peft}. La optimización utiliza AdamW con
$\eta = 1\times 10^{-4}$ y \emph{weight decay} $0{,}05$; la
pérdida combina \emph{cross-entropy} y \emph{Dice loss} sobre la
máscara binaria predicha; el tamaño de batch es $4$ (limitado
por el coste de SAM2.1-Large en BF16); la normalización es uniforme
con $\mu = \sigma = 0{,}5$ en los tres canales, según la
convención de SAM2. Se evalúan dos configuraciones: $50$ épocas
(resultado de referencia) y $100$ épocas con \emph{early
stopping} sobre el Dice de validación. El conjunto de entrenamiento
es ISIC2018~\cite{isic2018}; la transferencia fuera de dominio se
evalúa adicionalmente sobre ISIC2017 y PH2 sin reentrenamiento.
Como referencias se contrastan SAM2 \emph{zero-shot} sin adaptación
y un \emph{Convolutional AutoEncoder} (CAE-seg) entrenado desde
cero con la misma pérdida y configuración de optimización.

\subsection{Clasificación zero-shot multimodal}
\label{sec:zs}

El protocolo aprovecha el espacio de representación compartido
imagen-texto de los modelos vision-lenguaje contrastivos. Cada
clase candidata se formula como una o varias frases de plantilla;
las frases se codifican mediante el encoder textual del modelo y,
cuando existen varias plantillas por clase, los vectores resultantes
se promedian para obtener un único vector representativo. La imagen
se codifica con el encoder visual y la clase ganadora es la de
máxima similitud coseno con el vector visual. Sobre HAM10000 se
evalúan tres configuraciones de \emph{prompts}: (i)~genéricos en
inglés, una plantilla por clase
(p.~ej.~\emph{``a dermoscopic image of a melanoma''});
(ii)~conjunto A3 del paper de Derm1M~\cite{derm1m}, con varias
plantillas por clase y promediado posterior; (iii)~enriquecidos con
descriptores clínicos adicionales (modalidad, localización
anatómica, criterios diagnósticos visuales) sobre la base del
conjunto A3. El protocolo se aplica con dos tokenizadores distintos
para caracterizar la sensibilidad del rendimiento al tokenizador:
PubMedBERT extendido a $256$ tokens (DermLIP v1/v2) y GPT-2 con
tokenizador de CLIP a $77$ tokens (DermLIP original). Sobre HIBA,
MSKCC, DDI y Derm7pt clínico se ejecuta la versión binaria
melanoma/no-melanoma con \emph{prompts} genéricos y conjunto A3 de
Derm1M.

\subsection{Comparación con modelos de lenguaje multimodales}
\label{sec:llm}

El protocolo evalúa GPT-4o (API OpenAI), Gemini~2.5~Pro (API
Google), MedGemma~4B~Vision (local) y MedGemma~27B Text (local en
BF16) sobre el \emph{split} de test de HAM10000 ($1\,232$
imágenes), idéntico al utilizado en sondeo lineal y ajuste fino
para garantizar comparabilidad. El \emph{prompt} clínico
estandarizado de aproximadamente $300$ palabras describe la
tarea como apoyo al diagnóstico (no sustitución del juicio
clínico), presenta las siete categorías diagnósticas de HAM10000
con descripción breve y solicita el diagnóstico más probable con
razonamiento estructurado. La etiqueta predicha se extrae mediante
\emph{parsing} determinista que mapea el nombre del diagnóstico al
código correspondiente de HAM10000; las respuestas no mapeables
(menos del $1\%$ en todos los modelos) se contabilizan como
incorrectas. El modo principal es \emph{zero-shot} sin
\emph{fine-tuning} previo. Para MedGemma~27B se evalúa
adicionalmente una variante ajustada con LoRA ($r = 8$) sobre el
\emph{split} de entrenamiento de HAM10000.

\subsection{Evaluación de equidad por fototipo}
\label{sec:equidad}

El protocolo evalúa cinco encoders fundacionales (PanDerm Large,
PanDerm Base, DermLIP v2 en su rama visual, DINOv2 ViT-L/14 y
BiomedCLIP) sobre el dataset Fitzpatrick17k completo
($16\,577$ imágenes, $114$ patologías, $6$ fototipos
cutáneos I-VI). El protocolo común es sondeo lineal con la
configuración de secci\'on~\ref{sec:lp}, empleando para cada modelo su
normalización canónica. Se reportan dos formulaciones:
(i)~clasificación a $114$ patologías (grano fino del dataset)
y (ii)~clasificación a $3$ clases de malignidad
(benigno / no neoplásico / maligno), formulación habitual en
estudios de equidad sobre Fitzpatrick17k. Las métricas por subgrupo
se calculan sobre los seis fototipos. La métrica de equidad
declarada es la brecha entre los fototipos extremos
$\mathrm{Gap}_{\mathrm{I}\to\mathrm{VI}} =
\mathrm{BAcc}_{\mathrm{I}} - \mathrm{BAcc}_{\mathrm{VI}}$.
La verificación del experimento se realiza re-ejecutando el sondeo
sobre PanDerm Large y comprobando coincidencia hasta cuatro
decimales con el resultado original.

\subsection{Auditoría de solapamiento con Derm1M}
\label{sec:auditoria}

La auditoría caracteriza el solapamiento entre los \emph{splits} de
test contrastados y el dataset de preentrenamiento Derm1M. Cada
imagen se reduce a su hash MD5 mediante la biblioteca estándar de
Python; la intersección entre los conjuntos de hashes determina las
imágenes presentes simultáneamente en ambos datasets. La cardinalidad
se reporta por dataset en el capítulo~\ref{cap:resultados}. Los
doce datasets externos se auditan, con la excepción de SkinCon (no
evaluado sobre clases diagnósticas).

\section{Métricas}
\label{sec:metricas}

\paragraph{Clasificación.}
\emph{Balanced Accuracy} (BAcc, promedio de \emph{recall} por
clase) es la métrica principal en escenarios desbalanceados. La
\emph{accuracy} se reporta como referencia comparativa entre
modelos sobre la misma distribución de clases. La F1 ponderada se
calcula como media armónica de precisión y \emph{recall}
ponderada por la frecuencia de cada clase. La AUROC se calcula con
estrategia \emph{one-vs-rest} (\texttt{multi\_class='ovr'} en
\texttt{scikit-learn}) y promedio sobre clases. El coeficiente
kappa de Cohen corrige la concordancia por azar. \emph{Recall@$k$}
mide la presencia de la etiqueta verdadera entre las $k$ primeras
posiciones del ranking; se reporta específicamente para la clase
melanoma sobre HAM10000.

\paragraph{Segmentación.}
Coeficiente Dice
$\mathrm{DSC} = 2\,|P\cap G|/(|P|+|G|)$ e intersección sobre unión
$\mathrm{IoU} = |P\cap G|/|P\cup G|$, donde $P$ es la máscara
predicha y $G$ la de referencia.

\paragraph{Equidad.}
BAcc por subgrupo Fitzpatrick (I-VI) sobre la formulación de tres
clases de malignidad, y \emph{gap}
$\mathrm{Gap}_{\mathrm{I}\to\mathrm{VI}}
= \mathrm{BAcc}_{\mathrm{I}} - \mathrm{BAcc}_{\mathrm{VI}}$.
Valores próximos a cero indican rendimiento comparable entre
fototipos extremos; la interpretación requiere considerar la BAcc
global, ya que un \emph{gap} pequeño asociado a una BAcc también
pequeña reflejaría un modelo uniformemente inferior, no equitativo
en sentido sustantivo.

\paragraph{Operativas.}
Latencia por consulta en milisegundos (inferencia más posproceso)
y huella de memoria (VRAM en gigabytes en BF16), medidas sobre la
infraestructura descrita en secci\'on~\ref{sec:infraestructura}.

\section{Pruebas estadísticas}
\label{sec:estadistica}

Para los contrastes principales entre modelos sobre un mismo
conjunto de test se reportan intervalos de confianza al $95\%$
obtenidos mediante \emph{bootstrap} sobre $1\,000$ remuestreos del
conjunto de evaluación. Las comparaciones de AUROC entre dos
modelos sobre el mismo conjunto se evalúan mediante el test de
DeLong~\cite{delong1988}, que controla la correlación inducida por
la evaluación sobre el mismo conjunto. \textbf{[NO ESPECIFICADO:
confirmar si se aplicará corrección por comparaciones múltiples y
el método elegido (Bonferroni, Benjamini-Hochberg).]} Los tamaños
muestrales por subgrupo en el experimento de equidad se reportan
junto a las cifras correspondientes en el
capítulo~\ref{cap:resultados}.

\section{Trazabilidad y reproducibilidad}
\label{sec:reproducibilidad}

Las semillas aleatorias se fijan a $0$ para \texttt{torch},
\texttt{numpy} y \texttt{random} en los protocolos de
\emph{fine-tuning} y segmentación, y a $42$ en el parámetro
\texttt{random\_state} para los basados en \texttt{scikit-learn}
(sondeo lineal, eficiencia en etiquetas, equidad). El entorno de
ejecución queda fijado en las versiones de software declaradas en
secci\'on~\ref{sec:infraestructura}; cualquier desviación se documenta
explícitamente en el experimento correspondiente. La variable
\texttt{WANDB\_MODE} se fija a \texttt{disabled} para impedir el
envío de telemetría a servicios externos durante el
entrenamiento. La variable \texttt{CUDA\_VISIBLE\_DEVICES} se fija
a \texttt{0} para forzar la ejecución sobre la única GPU presente.
Los resultados numéricos se vuelcan a ficheros CSV o JSON con
esquema estable, lo que permite la reconstrucción de cualquier
tabla o figura del capítulo~\ref{cap:resultados} sin reejecución.
Los mapeos de etiquetas nativas a la ontología L1/L2/L3 se mantienen
fijos a través de los experimentos. Cuando un experimento emplea
como entrada el resultado de otro (por ejemplo, los \emph{embeddings}
extraídos en sondeo lineal alimentan la rama visual de zero-shot),
la identidad de la entrada se preserva mediante hash MD5 sobre el
fichero correspondiente. El listado completo de tareas
reproducibles, sus comandos de ejecución, ficheros de salida y
\emph{checkpoints} de modelo se documenta en el
Anexo~\ref{cap:anexoa}.
